\section{问题一模型建立与求解}
\label{sec:question-one}

问题一的决策从原计算图出发，依次确定计算操作的子图边界、各子图所在核心与核内先后。切分增加了可并行的Task，也改变张量读写次数、片上存活区间和Task激活时刻。为把这些作用放在同一口径下，本节先建立提交方案和执行约束，再由张量接收Task推导结构搬运，最后将有限候选送入官方事件评价，按返回结果选出每张图的方案。

\subsection{两级约束与评价目标}
\subsubsection{调度变量与执行约束}
第一层决定划分与次序：式\eqref{eq:plan-coverage}给出操作和子图覆盖，式\eqref{eq:plan-sequence-unique}要求每个非空子图在核内序列中恰好出现一次；若$s\in\sigma_k$，则$a(s)=k$。跨子图依赖形成的商图必须无环，同核序列还须遵守商图依赖；这两项检查在A、B、C三个场景共用。给定合法计划$X$后，评估程序展开Task与COPY，依次执行核内Step1排序、Step2换出恢复和Step3流水线及内存复用依赖构图，再模拟共享带宽事件。记这一展开为$\Psi_r(G,X)$，事件模拟为$\Phi_r$，其中$r\in\{A,B,C\}$。展开产生的操作集合记$\mathcal O_r(X)$，操作$o$的开始、结束时刻分别为$s_o,e_o$。数据或内存依赖$u\to v$要求$s_v\ge e_u$；带延迟的释放边要求$s_v\ge e_u+\delta_{uv}$。四条每核流水线M、V、MTE2和MTE3各自串行，不同流水线满足依赖后可重叠执行。

容量也属于第二层。逻辑张量在换出与恢复后可有不同物理版本$\tilde t$。若版本在核心$k$的存储区$q$上从$A_{\tilde t}$占用至$L_{\tilde t}$，任何时刻$\tau$均须满足
\begin{equation}
\sum_{\tilde t\in\mathcal T_{kq}} b_{\tilde t}
\mathbf1\{A_{\tilde t}\le\tau<L_{\tilde t}\}\le C_q,\qquad
C_{\rm L1}=524288,\quad C_{\rm UB}=131072\ \mathrm{B}.
\label{eq:capacity}
\end{equation}
这里的集合按物理版本而非逻辑张量计数。一个张量直到最后一个读者结束才释放，读者也包括向DDR写回的COPY。Step2在空间不足时生成换出恢复，Step3再为复用同一位置的操作补充先后约束。因此，改变Task边界不仅改变结构COPY，还可能改变物理版本的存活区间与恢复次数。

表\ref{tab:q1-hardware}列出执行递推所用的题设参数。原始文件的COPY操作周期是占位值；新生成的COPY独占参考工作量由字节数和带宽计算。问题一每个子图独立成Task，同核连续Task间隔100 cycle，跨核依赖Task释放延迟1000 cycle。后两问沿用容量与DDR参数，改变Task组织并在问题三加入Cache。
\begin{table}[H]\centering\caption{官方配置中的容量、带宽与Task等待}\label{tab:q1-hardware}
\begin{tabularx}{\linewidth}{lYY}\toprule
项目 & 数值 & 在执行递推中的作用\\\midrule
每核L1、UB & 524288 B、131072 B & 分存储区约束物理版本的同时驻留\\
共享DDR & 60 B/cycle & 各核COPY读写共同占用同一服务池\\
问题一同核等待 & 100 cycle & 同核前一Task结束后才能激活后一Task\\
问题一跨核等待 & 1000 cycle & 跨核前驱Task结束后延迟释放\\
问题二、三跨核延迟 & 500 cycle & 源端写出后目标COPY\_IN等待\\
问题三只读Cache & 1048576 B、250 B/cycle & 独立读池与FIFO容量，详见问题三\\\bottomrule
\end{tabularx}\end{table}

\subsubsection{共享DDR服务下的工期}
给定一次COPY的字节数$b_o$，其独占DDR参考工作量为$w_{o,D}=\max\{1,\lceil b_o/60\rceil\}$。若时刻$\tau$有$m_D(\tau)$个在途DDR搬运，各自剩余参考工作量$R_o(\tau)$按
\begin{equation}
R_o(s_o)=w_{o,D},\qquad
\frac{\mathrm dR_o(\tau)}{\mathrm d\tau}=-\frac{1}{m_D(\tau)}
\quad\bigl(o\text{正在DDR池服务}\bigr)
\label{eq:ddr-share}
\end{equation}
推进。一个搬运完成或新搬运发射后，$m_D$改变，其余搬运的预测完成时刻随之重算；随题官方多核模拟程序按整数周期规则处理事件。完整工期由该事件模拟确定：相同字节量的方案可能在关键路径上遇到不同的竞争，静态总字节与各COPY独占时间之和仅作量纲参考。

由此定义场景$r$的工期及额外搬运：
\begin{equation}
F_r(X)=\max_{o\in\mathcal O_r(X)}e_o,\qquad
D_r(X)=Q_r^{\rm sched}(X)-Q^{\rm orig}.
\label{eq:objective}
\end{equation}
以$\mathcal F_{r,K}$表示满足操作覆盖、商图无环、同核依赖顺序以及式\eqref{eq:capacity}物理容量约束的计划集合；Task等待和带宽服务按场景$r$由$\Psi_r$与$\Phi_r$确定。优化模型为
\begin{equation}
X^*_{r,K}\in\operatorname*{arg\,min}^{\rm lex}_{X\in\mathcal F_{r,K}}
\bigl(F_r(X),D_r(X)\bigr).
\label{eq:lex-objective}
\end{equation}
主目标是工期，工期相同时比较额外搬运；候选名仅用于程序的确定性平局处理。在仅有两核、每核一个Task、无依赖与搬运的特例中，各独立M操作的周期为正整数，判定工期能否不超过总周期的一半等价于将这些周期分成两个等和集合，即Partition判定。该特例已是弱NP完全问题，因此求解器生成有限候选，再以事件模拟评价可行计划。

\subsubsection{工期下界}
设$L_G$为计算依赖的加权最长路，$W_M,W_V$分别为两条计算流水线的总工作量；$Q_{\min}$为图输入至少一次读入与最终输出至少一次写出的最低字节数，$B_D=60$~B/cycle为共享DDR服务带宽。忽略Task等待、容量恢复、流水线排序和争用，先分别得到计算侧和IO侧下界
\begin{equation}
LB_{\rm comp}(K)=\max\left\{L_G,\frac{W_M}{K},\frac{W_V}{K}\right\},
\qquad
LB_{\rm io}=\frac{Q_{\min}}{B_D}.
\label{eq:lower-bound-components}
\end{equation}
综合松弛下界取两侧瓶颈中的较大者：
\begin{equation}
LB(K)=\max\left\{LB_{\rm comp}(K),LB_{\rm io}\right\}.
\label{eq:lower-bound}
\end{equation}
其中$LB_{\rm comp}$由依赖串行段和两条计算流水线的吞吐约束共同给出，$LB_{\rm io}$只刻画最低DDR服务量；$LB$保留两侧中较大的不可绕过瓶颈。该下界由依赖最长路、M/V工作量和最低DDR服务量构成，未计入Task等待、容量恢复、流水线排序与共享带宽争用；$F/LB$用于诊断方案与理想瓶颈的距离。各场景加速比使用同图整图单Task、单核A0的官方工期$T_{1,i}$作为统一速度比基准，使问题二、三的收益与问题一比较时具有一致起点。

\subsection{Task驻留域与结构搬运}
\subsubsection{Task划分与边界搬运}
问题一中每个子图形成独立Task。Task内部可以让同一张量服务多个计算消费者；跨Task时，即使两个Task安排在同一核心，后一个Task也要经过DDR边界重新读入。结构搬运按接收Task数量计量，跨核依赖边数不足以给出完整通信量。对输入张量$t$，记$I_A(t)=\{g(v):v\in U(t)\}$；对计算产生的张量，记$R_A(t)=\{g(v):v\in U(t),\ g(v)\ne g(p(t))\}$。集合消除了同一Task内的重复消费者。由官方Task构图规则可逐张量列账：输入对每个接收Task读入一次；内部结果若有外部接收Task，源Task写出一次，各目标Task各读入一次；最终输出仍需源Task写出。因此，不含容量恢复的结构COPY字节数为
\begin{equation}
Q_A^0=\sum_{t\in\mathcal T_{\rm in}}b_t|I_A(t)|
+\sum_{t\in\mathcal T_{\rm prod}}b_t
\left[|R_A(t)|+\mathbf1\{|R_A(t)|>0\text{ 或 }o_t=1\}\right].
\label{eq:a-bytes}
\end{equation}
例如，大小为$b$的内部张量由一个Task生产，在另外两个Task各有一个消费者，且不是最终输出。其结构搬运为一次写出与两次读入，共$3b$；将两个消费Task合并后变为$2b$；连同生产者一起合并则这部分边界搬运为零。若只把两个消费Task放到同一核心而不合并，问题一仍保留$3b$。这说明计算均衡和边界数并非同一目标：切开一条并行支路时，应同时检查边界附近的大张量。

实际调度COPY还包含容量不足时的换出与恢复。将官方字段对应到模型，得到
\begin{equation}
Q_A^{\rm sched}=Q_A^0+Q_A^{\rm spill},\qquad
D_A=\underbrace{(Q_A^0-Q^{\rm orig})}_{\text{划分新增搬运}}
+\underbrace{Q_A^{\rm spill}}_{\text{容量恢复搬运}}.
\label{eq:a-copy-breakdown}
\end{equation}
式\eqref{eq:a-bytes}计算切图决定的逻辑边界，式\eqref{eq:a-copy-breakdown}将其与片上容量造成的物理搬运分开。即使结构搬运降低，张量生产和末次使用之间的间隔变长，也可能在Step2产生更多恢复字节；反之，适度切分虽增加边界COPY，却可能缩短张量的单Task存活区间。最终是否加速取决于这些COPY落在何时和哪条依赖路径上。

\subsubsection{Task激活时刻}
令Task $s$的激活、完成时刻为$B_s,E_s$，同核前一个Task为$\operatorname{prev}(s)$，其计算数据前驱Task集合为$\operatorname{Pred}(s)$。激活需同时等待本核前一Task、所有数据前驱和跨核释放延迟：
\begin{equation}
B_s\ge\max\left\{0,\ E_{\operatorname{prev}(s)}+100,\
\max_{t\in\operatorname{Pred}(s)}E_t,\
\max_{\substack{t\in\operatorname{Pred}(s)\\a(t)\ne a(s)}}(E_t+1000)\right\},
\label{eq:a-wait}
\end{equation}
不存在的前驱项省略；各核心第一个无前驱Task可在零时刻激活。式中取最大值而非求和：若同核等待和数据依赖同时发生，较晚的一项决定可激活时刻。切分使不同核心出现并行机会，也可能把一个原本位于Task内部的依赖改成相隔1000 cycle的跨核释放。该等待与结构COPY在同一切点出现，故应把二者连同计算负载一起比较。

切块程序沿拓扑序累计M、V两类计算周期，以两条流水线工作量的较大者接近目标值作为负载条件，并在接近目标的切点中优先减少跨切点张量字节。生成完整计划后，轻量分数再计入路径等待、结构搬运与容量压力；工期取式\eqref{eq:objective}的事件模拟结果。若同一张量被多个新Task消费，式\eqref{eq:a-bytes}中的$|R_A(t)|$随接收Task数增加。

\subsection{方案构造与求解}
\subsubsection{初始候选方案}
单个大Task保留数据驻留但难以让多核并行；只使用弱连通分量又不能切开\Case{016}一类近乎单分量的大图。程序以四个确定性主起点提供整图、天然独立支路和两种拓扑切分，并保留四个按操作数切块的旧起点供展开阶段比较。表\ref{tab:a-seeds}列出主起点。
\begin{table}[H]\centering\caption{问题一的四种确定性候选起点}\label{tab:a-seeds}
\begin{tabularx}{\linewidth}{cYY}\toprule
名称 & 子图构造 & 对应的调度取舍\\\midrule
A0 & 整图为一个子图，置于核心0 & 保留完整Task驻留，作为独立单核基准计划\\
A1 & 弱连通分量按M/V两类工作量装入$K$个桶，每个非空桶形成一块 & 利用天然并行并平衡两条计算流水线\\
A2 & D序按M/V工作量接近目标、跨切点字节较少的规则切至多$2K$块 & 为大分量增加切点，维持分量内先后\\
A3 & H序按相同规则切至多$4K$块 & 让关键链附近产生另一组边界和核内次序\\\bottomrule
\end{tabularx}\end{table}

H序的优先级由逆拓扑递推
\begin{equation}
h(v)=w_v+\max_{u\in\operatorname{Succ}(v)}h(u),\qquad
\max\varnothing=0
\label{eq:a-height}
\end{equation}
得到。只有前驱已放入序列的操作才参与就绪队列；$h(v)$相同时按编号确定次序。D序也遵守拓扑顺序，但优先保持同一弱连通分量的操作相邻。这两种次序都只是生成切点的输入，不是最终硬件的执行时间线。A2、A3的块数上限随$K$增加。对已安置前$j-1$块的部分计划$X^{(j-1)}$，记$X^{(j-1)}\oplus_k S_j$为把第$j$块接到核心$k$当前序列末尾后的部分计划。逐块比较轻量评分，选择
\begin{equation}
a(S_j)=\operatorname*{arg\,min}_{0\le k<K}
\left(\widehat\kappa_A\bigl(X^{(j-1)}\oplus_k S_j\bigr),k\right).
\label{eq:a-place}
\end{equation}
编号作为评分相同时的确定性次序。式中的$\widehat\kappa_A$是部分计划的轻量评分，包含路径、计算负载、结构COPY与容量压力；其值不等于实际工期。A1先按M/V工作量装桶，再以每桶一个子图形成方案；块内操作由Step1确定实际顺序。

\subsubsection{候选方案调整}
四起点各自保留原方案，并生成迁块、同核合并、调序及容量风险块前后移动等变体。先对起点做官方profile，以展开后的排序量选一个搜索起点；再在有限轮局部搜索中考察子图迁移、拆分与合并、边界操作迁移和核内重排，保留轻量分数改善代表及容量压力下降代表。改动后的映射与列表先规范化，检查覆盖、商图无环及同核依赖顺序，再按计划内容去重；非法或重复变体跳过。

上一核数的已选计划在新增核心保持空闲后，作为本核数的继承候选。整图、分量装桶和继承计划按内容去重后分别取得必评角色；同一计划承载多个角色时只评价一次。文中的“初始到最终”比较表示已评候选之间的选择差异，用于衡量搜索收益。

\subsubsection{官方评价与方案选定}
TVM以学习型成本模型加快低层程序优化空间的探索\cite{chen2018tvm}，CoSA则将空间加速器的时空映射纳入约束优化\cite{huang2021cosa}。本题的完整模拟会为每个Task生成COPY、物理版本和事件，因而先用静态评分排列有限候选。设$\widehat L(X)$为沿计算拓扑序的路径代理；跨子图依赖在该代理中加入固定等待权重。设$\widehat Q(X)$为按张量接收域估计的COPY字节，$\widehat\Pi(X)$为每核按局部拓扑操作序估计的L1/UB超容量峰值。程序的第一排序量是
\begin{equation}
\widehat F(X)=\max\left\{\widehat L(X),\ \max_k L_k,\
\left\lceil\frac{\widehat Q(X)}{60}\right\rceil\right\},\qquad
\widehat\kappa(X)=\bigl(\widehat F(X),\widehat Q(X),\widehat\Pi(X)\bigr).
\label{eq:a-light-score}
\end{equation}
其中$\widehat Q$和$\widehat\Pi$只用于安排试评顺序：前者按接收域估计结构搬运，后者以每核操作序上的张量首次与末次使用估计容量压力，均不代替官方展开。路径代理对跨子图依赖加1000 cycle，未细分同核Task的100 cycle等待，也未模拟流水线与共享DDR事件；最终报告采用完整事件评价返回的$F_A$与$D_A$。式\eqref{eq:a-place}的$\widehat\kappa_A$采用同一排序量在部分计划上的取值。

官方完整评价先覆盖去重后的整图、分量装桶和上一核继承计划，再从其余候选中至少评价一个代表。设去重后必评计划数为$m$、请求额度为$q$，常规完整调用额度为$\max\{\min(q,2),m+1\}$。待评方案先调用官方profile；展开失败保留错误并转向下一份，不占完整调用额度，完整评价失败则占用一次并记录错误。成功集合按$(F_A,D_A,\text{候选名})$选优；对目标核数$K$，再将$1,\ldots,K$核的成功方案与整图基准逐图比较，取Makespan最短者并以空核列表补齐。正文曲线、逐图长表和慢例统计均按该继承规则汇总。

\begin{algorithm}[H]\caption{问题一有限候选生成与官方评价}
\begin{algorithmic}[1]
\REQUIRE 原始图$G$、核数$K$、官方参数及请求完整调用额度$q$
\ENSURE 选中映射、核心列表、工期和额外搬运
\STATE 预处理生产消费关系，计算拓扑序、分量、D序和H序
\STATE 生成A0整图、A1分量装桶、A2的$2K$块与A3的$4K$块
\FOR{每个起点}
\STATE 保留原方案；生成迁块、合并、调序和容量风险变体
\STATE 规范化，检查商图及同核顺序，按方案内容去重
\ENDFOR
\STATE profile各起点并搜索轻量改善与容量改善代表
\STATE 置整图、分量装桶和上一核继承计划为去重后的必评计划
\STATE 令$m$为必评数，完整调用额度为$\max\{\min(q,2),m+1\}$
\WHILE{尚有候选且未达到常规完整调用额度}
\STATE 先取未评必评计划，再取已profile代表或轻量评分靠前者
\STATE profile失败则记录并继续；完整事件模拟保存结果或失败文本
\ENDWHILE
\STATE 若当前最优慢于上一核或单核基准，补评尚未评价的回退计划
\STATE 在成功集合按$(F_A,D_A,\text{候选名})$选最终方案
\end{algorithmic}\end{algorithm}

伪代码已经列出必评锚点、探索额度和回退关系。静态分数及profile排序只决定试评顺序，工期、搬运和容量结果来自官方事件评价；两类失败分别留痕。整图单核A0工期是统一速度比基准，A、B单核搜索结果另作候选诊断。

表\ref{tab:a-code-map}列出与模型步骤对应的实现模块和输出对象；\texttt{profile.json}记录展开状态，\texttt{result.json}记录完整模拟工期。
\begin{table}[H]\centering\small\caption{问题一实现模块与模型输出}\label{tab:a-code-map}
\begin{tabularx}{\linewidth}{lYY}\toprule
实现模块 & 模型步骤 & 输出对象\\\midrule
\texttt{graph.py}、\texttt{plan.py} & 生产消费索引、D/H序、四起点与计划合法性检查 & 依赖索引、映射和核内序列\\
\texttt{candidates.py}、\texttt{scoring.py} & 局部变体、去重与轻量排序 & 候选方案和评分向量\\
\texttt{official.py}、官方A评估器 & Task展开、profile与共享DDR事件模拟 & Task/COPY展开、工期和搬运\\
\texttt{official.py}的选优接口 & 成功结果按工期、搬运和候选名排序 & 初始方案与最终方案\\\bottomrule
\end{tabularx}\end{table}

\subsection{结果口径与分析}
\subsubsection{结果指标与速度比}
三问使用同一张图的整图单Task、单核A0官方工期$T_{1,i}$作基准。问题一正式需要2至5核各100张图，共400组；另保存100组单核搜索结果以诊断候选行为。按照题面单核曲线点的口径，令
\begin{equation}
S_{A,i,K}=\begin{cases}
1,&K=1\text{ 的正式曲线点},\\
T_{1,i}/F_A(X^*_{i,K}),&K=2,3,4,5,
\end{cases}\qquad
\overline S_{A,K}=\frac1{100}\sum_{i=1}^{100}S_{A,i,K}.
\label{eq:a-speedup}
\end{equation}
这是先对每张图求工期比、再对100个比值取算术平均，而非用100张图的基准工期总和除以最终工期总和。对目标核数$K$，在已评价的$1,\ldots,K$核方案和整图单核基准中取最短Makespan；低核方案以空核列表补齐后作为$K$核方案。各核数均使用同一批100张图，失败候选不参加比值的分子或分母。

表\ref{tab:a-results}给出正式曲线均值和分布信息。图\ref{fig:speedup-a}的浅色带对固定100张图作1500次有放回实例重采样，取均值的2.5\%和97.5\%经验分位数，反映图集合组成改变时的均值波动，不代表芯片重复运行噪声。
\begin{table}[H]\centering\small\caption{问题一：100图逐核工期比与候选改善（慢例按未舍入工期严格比较）}\label{tab:a-results}
\begin{tabularx}{\linewidth}{cYYYYY}\toprule
核数 & 平均加速比 & 中位数 & 均值实例重采样区间 & 慢于A0图数 & 原候选选择收益/\%\\\midrule
1 & 1.0000 & 1.0000 & [1.0000,1.0000] & 0 & 0.8800\\
2 & 1.6876 & 1.9626 & [1.6003,1.7695] & 0 & 5.0955\\
3 & 2.3176 & 2.8340 & [2.1538,2.4716] & 0 & 8.5596\\
4 & 2.9316 & 3.6268 & [2.6855,3.1642] & 0 & 11.6031\\
5 & 3.4132 & 3.9605 & [3.1025,3.7151] & 0 & 16.9668\\\bottomrule
\end{tabularx}\end{table}
\samplefigure{speedup_A}{问题一加速比：100图均值、中位数及均值95\%实例重采样区间}{fig:speedup-a}

五核继承后平均加速比为3.4132，中位数为3.9605。均值随核数从1、1.6876、2.3176、2.9316增至3.4132，逐段增加0.6876、0.6300、0.6140、0.4816。4至5核的61张图采用更短的已评价计划，39张图沿用四核或整图基准计划，因此没有逐图退化；表中最后一列仍保留原候选集合的选择收益，便于区分继承规则和候选搜索。

\subsubsection{搬运量分解}
表\ref{tab:a-copy-results}把额外搬运按式\eqref{eq:a-copy-breakdown}拆开。每格是先对100张图逐图得到官方字节数，再转换为MiB求算术平均；一个图的额外搬运可以因大图而远高于全体中位数。随着核数由1增至5，平均划分新增搬运从0.275升至4.001 MiB，平均容量恢复则从11.348降至9.214 MiB。二者方向相反，故总额外搬运并不随核数单调变化。五核有28张图发生容量恢复，说明即使平均恢复量下降，片上空间仍是部分图的显著成本。
\begin{table}[H]\centering\caption{问题一额外搬运的结构与容量分解（100图均值，MiB）}\label{tab:a-copy-results}
\begin{tabular}{crrrr}\toprule
核数 & 划分新增 & 容量恢复 & 额外搬运合计 & 有容量恢复的图数\\\midrule
1 & 0.275 & 11.348 & 11.624 & 37\\
2 & 1.292 & 10.365 & 11.657 & 35\\
3 & 2.069 & 10.648 & 12.717 & 34\\
4 & 3.022 & 10.388 & 13.409 & 31\\
5 & 4.001 & 9.214 & 13.216 & 28\\\bottomrule
\end{tabular}\end{table}

由此不能用“COPY越少，工期越短”直接选解。式\eqref{eq:ddr-share}说明同样的字节量可以在不同的时间遇到不同并发数；式\eqref{eq:a-wait}又使跨Task和跨核的依赖等待影响关键路径。候选选择阶段以字节量辅助排序，最终仍以模拟工期为第一目标。这也解释了为什么表\ref{tab:a-copy-results}中五核平均额外搬运略低于四核，却不能保证每张图的五核工期都更短。

\subsubsection{典型与退化算例}
\Case{001}有600个非COPY操作。其五核选中A1方案把它们分为五个120操作的子图，\texttt{core\_schedules}为\texttt{[[0],[1],[2],[3],[4]]}；每核执行一块。映射中操作3、4、5和83、84、85均属于子图0，完整字段覆盖600个操作。官方整图单核工期233110 cycle，四核58984 cycle，五核47502 cycle，五核加速比4.9074。该五核方案新增划分搬运4608 B，容量恢复为零。这里的并行收益与极小的新增搬运同时出现；实际工期仍包含Task激活和DDR事件，不能把4.9074写成纯计算负载比。

\Case{053}展示容量约束持续起作用的情形。其五核选中方案工期492699 cycle，对整图单核1472686 cycle的比值为2.9890；官方额外搬运6918912 B，其中结构新增3104256 B、容量恢复3814656 B。四核的492232 cycle反而比五核快467 cycle。这里仅能从字节账本确认容量恢复存在并占额外搬运的一部分；是否位于决定467 cycle差异的关键路径，需要对两份时间线逐事件比较。

另外两张图显示核数退化并不只有一种字节形态。\Case{055}的四核工期32099 cycle，五核110002 cycle；五核方案为10个子图，额外搬运199780 B且没有容量恢复。\Case{067}则从四核18473155 cycle升至五核63959635 cycle，五核额外搬运193700096 B，其中容量恢复167278592 B；它的五核容量恢复还低于四核的252010112 B。因此，单独减少恢复字节并不足以保证工期下降。两例都要求把被选中的Task顺序、释放等待和并发DDR时间线与搬运量一起看，而不能用一个平均字节数给全部慢例归因。

\subsubsection{核数变化下的依赖结构}
应用继承规则后，4至5核的100图配对中61图采用更短工期、39图工期不变，没有图变慢；五核方案均不低于整图单核基准。绝对工期均值随逐图取最短方案下降，四核为1035607.10 cycle，五核为921892.21 cycle。该结果将核数选择转化为逐图可行的方案选择，均值和单图工期均采用同一继承口径。

五核均值的实例重采样区间为[3.1025,3.7151]，两核区间为[1.6003,1.7695]。重采样以图为单位保留每张图内部的基准与继承后工期配对，检验结论对图集合组成的敏感度。容量、带宽变体的专门实验置于问题三。

\subsubsection{方案合法性检查}
每个最终方案回代\texttt{node\_to\_subgraph}和\texttt{core\_schedules}，检查全部非COPY操作恰好归属一次、子图不空且在核心列表中恰好出现一次，再由官方入口展开执行图。输出同时保存候选名、计划内容、计划哈希、原始与调度COPY、容量恢复、时间线及最晚完成时刻。结果包对1500个已存组合键和其中2784条成功候选记录做了数值、计划和时间线一致性检查；A正式部分为2至5核的400组，单核100组作为辅助诊断。另有包含A场景在内的15项跨场景官方抽样重算，与对应存档结果逐项相同。结果文件的字段一致性检查与抽样重算方法见复现附录。

问题一的交付是每张图、每个核数的具体子图映射及核内列表。当前有限候选中，五核有68图选A1家族、26图选A2家族、3图选A3家族、3图选A0家族；39图的最终工期低于首个成功候选，逐图相对降幅的均值为16.9668\%。这说明保留多种切分起点确实改变了入围方案，但“首个成功到最终”的比较不是移除某个机制的独立消融。正式工期来自成功完成官方评价的计划，未评价过的划分没有进入最优性判断。逐图工期、额外搬运及五核慢例可在附录\ref{app:results}回查。

\subsection{问题一灵敏度分析}
以可用核数$K$为扰动量，固定100张原图和场景A规则，在$1,\ldots,K$核已评价方案与整图基准中逐图取最短工期。由四核增至五核时，61图缩短、39图不变、0图延长，平均加速比由2.9316升至3.4132。图\ref{fig:speedup-a}给出总体曲线，表\ref{tab:a-results}给出继承后的零慢例结果。

\Case{067}的四核与五核工期分别为18473155和63959635 cycle，增加一个核心后工期变为原来的3.462倍；其额外搬运量却从263174912 B降至193700096 B。该反例表明，边界字节数的变化不足以预测完工时刻，Task等待与共享DDR的事件位置也会改变尾部工期。对这类图，实际选核应比较已经通过官方评价的完整计划，优先保留较少核心下更短的方案。
